\section{Theory}
\label{sec:theory}

\subsection{Reversible chains and the discriminant matrix}
\label{sec:theory-discriminant}

Let $P$ be an ergodic reversible Markov chain with stationary distribution
$\pi$. The discriminant matrix is
\begin{equation}
  D_{xy}=\sqrt{P_{xy}P_{yx}},
  \qquad
  D=\mathrm{diag}(\pi)^{1/2}P\,\mathrm{diag}(\pi)^{-1/2}.
  \label{eq:discriminant}
\end{equation}
For reversible $P$, $D$ is symmetric and has the same spectrum as $P$; the
vector $\sqrt{\pi}$ is a left eigenvector with eigenvalue $1$. The Szegedy walk
\eqref{eq:background-szegedy} satisfies
\begin{equation}
  \langle\bar0|\langle x|W(P)|\bar0\rangle|y\rangle=D_{xy}.
\end{equation}
The following standard spectral fact is used repeatedly. If
$D=\sum_k\lambda_k|v_k\rangle\langle v_k|$, then the nontrivial eigenvalues of
$W(P)$ are $e^{\pm i\theta_k}$ with $\lambda_k=\cos\theta_k$, and the
corresponding eigenvectors lie in the span of
$\{\lvert\bar0\rangle\lvert x\rangle\}_{x\in V}$ and its image under $V(P)$.
This is the Szegedy spectral lemma; we use it as a statement and do not
reprove it.

\subsection{Classical and extended hitting times}
\label{sec:theory-ht}

For a marked set $M$ with stationary mass $p_M=\sum_{x\in M}\pi_x$, the
classical hitting time can be written in spectral form as
\begin{equation}
  \HT(P,M)=\frac{1}{1-p_M}\sum_{k=1}^{n-|M|}
  \frac{\lvert\langle v'_k\rvert\sqrt{\pi_U}\rangle\rvert^2}{1-\lambda'_k},
  \label{eq:ht-spectral}
\end{equation}
where $\lambda'_k$ are the non-unit eigenvalues of $D(1)$ and
$\sqrt{\pi_U}$ is the subnormalized vector with entries $\sqrt{\pi_x}$ for
$x\notin M$.

AGJK replace $D$ by the discriminant matrix $D(s)$ of the interpolated walk and
define
\begin{equation}
  \HT(s)=\frac{1}{1-p_M}\sum_{k=1}^{n-1}
  \frac{\lvert\langle v_k(s)\rvert\sqrt{\pi_U}\rangle\rvert^2}{1-\lambda_k(s)},
  \qquad
  \HTp=\lim_{s\to1}\HT(s),
  \label{eq:ht-plus-theory}
\end{equation}
where $D(s)=\sum_k\lambda_k(s)|v_k(s)\rangle\langle v_k(s)|$ and
$\lambda_n(s)=1$. For a single marked vertex, $\HTp=\HT$. For multiple marked
vertices, \eqref{eq:ht-plus-theory} can be much larger than \eqref{eq:ht-spectral}.

The canonical example is an $N\times N$ torus with marked set
\begin{equation}
  M=M_1\cup M_2,\qquad
  M_1=\{(j_1d_1,j_2d_1):0\le j_1,j_2<k_1\},\quad
  M_2=\{(j_1d,j_2d):0\le j_1,j_2<N/d\},
  \label{eq:adversarial-set}
\end{equation}
with $d_1\mid d\mid N$ and the scale separation of AGJK. The set $M_1$ is a
dense cluster and $M_2$ is a sparse periodic grid. For suitable parameters,
$\HT=\widetilde O(d^2)$ while $\HTp=\Omega(N^2)$; this is AGJK's asymptotic
statement. Our numerics ($L=18$--$36$) can show growth of $\HTp/\HT$ but cannot
establish the asymptotic rate. The large gap comes from the
periodic structure, not from the mere fact that $|M|>1$.

\subsection{Interpolated walk and the numerical AGJK protocol}
\label{sec:theory-interpolated}

Let $U=V\setminus M$. The absorbing and interpolated walks are
\begin{equation}
  P'=\begin{pmatrix}P_{UU}&P_{UM}\\0&I\end{pmatrix},
  \qquad
  P(s)=(1-s)P+sP',\quad s\in[0,1).
  \label{eq:interpolated-theory}
\end{equation}
The quantum version attaches a qubit that records whether the current vertex is
marked. For $x\in U$, $V(P,s)$ acts as $V(P)$; for $x\in M$ it maps
\begin{equation}
  \lvert0\rangle\lvert\bar0\rangle\lvert x\rangle
  \mapsto
  \sqrt{1-s}\lvert0\rangle V(P)\lvert\bar0\rangle\lvert x\rangle
  +\sqrt{s}\lvert1\rangle\lvert\bar0\rangle\lvert x\rangle.
\end{equation}
The interpolated walk operator is
\begin{equation}
  W(s)=V^\dagger(P,s)\,\mathrm{Shift}'\,V(P,s)\,\mathrm{Ref}',
  \label{eq:interpolated-walk}
\end{equation}
with the controlled shift and reflection of AGJK. Its top-left block is
$D(s)$.

Algorithm~1 of AGJK prepares $\lvert\bar0\rangle\lvert\sqrt{\pi}\rangle$,
applies $W(s)$ for $t$ steps, and measures the marking ancilla. If the ancilla
is $1$, a position measurement returns a candidate marked vertex. The key
numerical object is the lower bound
\begin{equation}
  q_t(s)=\left\lVert\Pi_M T_t(D(s))\lvert\sqrt{\pi}\rangle\right\rVert^2
  \le
  \left\lVert(I\otimes\Pi_M)W^t(s)\lvert\bar0\rangle\lvert\sqrt{\pi}\rangle
  \right\rVert^2,
  \label{eq:chebyshev}
\end{equation}
where $T_t$ is the Chebyshev polynomial of the first kind. The identity behind
\eqref{eq:chebyshev} is
\begin{equation}
  \Pi_0W^t(s)\Pi_0
  =\lvert\bar0\rangle\langle\bar0\rvert\otimes T_t(D(s)),
  \qquad
  \Pi_0=\lvert\bar0\rangle\langle\bar0\rvert\otimes I.
\end{equation}
In our numerical protocol we scan $s$ over a fixed grid and choose the value
that maximizes $\max_t q_t(s)$ over $t\le3\sqrt{\HT}+5$. We then run the exact
operator $W(s)$ and record the actual success curve $p_M(t)$. This tests the
numerical content of Conjecture~3 of AGJK without assuming it.

For completeness, the proved AGJK algorithm combines the interpolated walk with
quantum fast-forwarding and runs in time
\begin{equation}
  O\!\left(S\sqrt{\log\HT}
  +\sqrt{\HT}(U+C)\sqrt{\log\HT\log\log\HT}\right),
  \label{eq:agjk-theorem}
\end{equation}
where $S$, $U$, and $C$ are the costs of setup, update, and checking. We cite
\eqref{eq:agjk-theorem} as the theoretical guarantee; our experiments use the
simpler Algorithm~1.

\subsection{Partial measurement}
\label{sec:theory-measurement}

The decoherence model is a position measurement with probability $p$ before
each unitary step. Let
$\Pi_x=I\otimes\lvert x\rangle\langle x\rvert$ be the projector onto vertex
$x$ and let $R_x=\lvert\bar0\rangle\lvert x\rangle$ be the reset state. The
one-step channel is
\begin{equation}
  \rho_{t+1}
  =W(0)\left[(1-p)\rho_t+p\sum_xR_x\langle R_x|\rho_t|R_x\rangle\langle R_x|\right]W(0)^\dagger .
  \label{eq:measurement-channel}
\end{equation}
For $p=0$ this is the pure quantum walk. For $p=1$ the state is reset to a
position basis state at every step; the position distribution then follows the
classical Markov chain
\begin{equation}
  Q_{xy}=\left\lVert\Pi_yW(0)\lvert\bar0\rangle\lvert x\rangle\right\rVert^2.
\end{equation}
For $s=0$, $Q=P$ and the measured walk is exactly the classical random walk.

Because \eqref{eq:measurement-channel} is linear in $\rho$, the average success
probability can be computed without Monte Carlo. We precompute the pure states
$\lvert\phi_{x,a}\rangle=W(0)^a\lvert\bar0\rangle\lvert x\rangle$ and their
position marginals and marked masses. A trajectory is summarized by the
probability $q_0$ that no measurement has occurred and by the distribution
$q_{x,a}$ over the last reset vertex $x$ and age $a$. The update is a small
linear recursion; the success probability is the corresponding mixture of the
precomputed marked masses. The details are in \cref{app:measurement}.

\subsection{Continuous-time walks on semantic graphs}
\label{sec:theory-ctqw}

For the semantic experiment we use the symmetric normalized Laplacian
\begin{equation}
  L_{\mathrm{sym}}=I-D^{-1/2}AD^{-1/2}
\end{equation}
and the random-walk generator
$L_{\mathrm{rw}}=I-D^{-1}A$. The CTQW and CRW are
\begin{equation}
  \lvert\psi(t)\rangle=e^{-iL_{\mathrm{sym}}t}\lvert c\rangle,
  \qquad
  \rho(t)=\rho(0)e^{-L_{\mathrm{rw}}t},
  \label{eq:ctqw-crw-theory}
\end{equation}
where $c$ is the cue vertex. The eigenvalues of $L_{\mathrm{sym}}$ and
$L_{\mathrm{rw}}$ coincide, so comparing at $\tau=\lambda_1t$ aligns the two
relaxation scales. The target activation is
$p_T(t)=\lvert\langle T\vert\psi(t)\rangle\rvert^2$ for CTQW and
$p_T(t)=\rho_T(t)$ for CRW.
